% sec:rk-fsal — First Same As Last (FSAL)
\section{First Same As Last (FSAL)}
\label{sec:rk-fsal}

Every accepted step of the adaptive integrator calls the right-hand
side $f$ seven times --- once per stage.  For the geodesic integrator,
each call computes the Hamiltonian right-hand side
$(\dot{x}^\mu, \dot{p}_\mu)$ from the current phase-space state, which
involves evaluating the metric, its inverse, and its derivatives at
the stage point.  Over a typical geodesic trace of $10^4$ steps, that
is $7 \times 10^4$ evaluations.  If the tableau structure let us reuse
one evaluation from the previous step, we would cut the per-step cost
by a seventh for free.  The FSAL property does exactly this.

\paragraph{The structural coincidence.}

Recall from \cref{eq:dp-tableau} that the primary weights
$\mathbf{b}$ --- which combine the stage derivatives into the
fifth-order solution --- are identical to the last row of the coupling
matrix $A$:
%
\begin{equation}
  b_i = a_{7,i} \,, \qquad i = 1, \ldots, 6 \,.
  \label{eq:fsal-condition}
\end{equation}
%
This means that the fifth-order solution
$y_{n+1} = y_n + h\sum_{i=1}^{6} b_i\,k_i$ (recall $b_7 = 0$) is
precisely the state at which the seventh stage is evaluated:
$k_7 = f(t_{n+1},\, y_{n+1})$.  But the first stage of the
\emph{next} step also evaluates $f$ at the accepted state:
$k_1^{(n+1)} = f(t_{n+1},\, y_{n+1})$, where the superscript labels
the step.  These are the same function call --- the last stage of
step $n$ is the first stage of step $n+1$:
%
\begin{equation}
  k_7^{(n)} = k_1^{(n+1)} \,.
  \label{eq:fsal-reuse}
\end{equation}
%
The name captures the identity: the \emph{first} stage of the new step
is the \emph{same as} the \emph{last} stage of the preceding one.

\paragraph{Cost savings.}

The first step of an integration has no predecessor, so all seven stages
must be computed from scratch.  Every subsequent accepted step reuses
$k_7$ from the previous step as $k_1$, computing only stages
$k_2$ through $k_7$ --- six new evaluations rather than seven.  Over
a long integration the FSAL property eliminates one right-hand-side
evaluation per accepted step, reducing the effective per-step cost from
seven to six function calls.
\implnote{The \texttt{dormandPrinceStep} function in
\codemodule{Geodesic.Solver} accepts $k_1$ as an argument and returns
the new state, the error estimate, and $k_7$ as a triple.  The
adaptive loop passes $k_7$ from each accepted step as $k_1$ of the
next; the initial $k_1$ is computed once at the start of integration.}

When a step is rejected (because the error estimate exceeds the
tolerance), the integrator retries from the same state with a smaller
step size.  The existing $k_1$ remains valid --- it was computed at the
current state, which has not changed --- so no additional evaluation is
wasted.

\paragraph{Design constraint.}

The FSAL property is not automatic.  It requires the tableau to be
designed so that \cref{eq:fsal-condition} holds, which consumes some
of the free parameters that would otherwise be available for
minimising truncation-error coefficients.  Dormand and Prince
incorporated this constraint into their optimisation
\cite{DormandPrince1980}, accepting a slightly larger truncation-error
coefficient in exchange for the guaranteed one-evaluation saving per
step --- a worthwhile trade for any method that will be used with
adaptive step-size control, where the total number of steps is large.
For the geodesic integrator, the net effect is simple: the first step
pays full price, and every subsequent accepted step gets one evaluation
free --- a saving that compounds across millions of rays.
